\section{Evaluation of \sys{}}\label{sec:eval}

\subsection{What a name-level profile shows}\label{sec:namelevel}

Profilers such as Nsight Systems summarize kernels by name. \Cref{tab:namelevel} applies that
grouping and \sys{}'s grouping to the same 38\,880 launches of the batch-1 capture. By name, the
CUTLASS FP8 GEMM is a single row: 148 launches and 66.7\,ms per step, spread over seven weight
groups, with a coefficient of variation of 1.92. Nothing in that row distinguishes a 47\,\textmu s
excess on one weight group from the expected spread of different shapes. Split by weight, shape
and grid, each group is tight (the 10th to 90th percentile range is under 4\% of the median), and the GDN \texttt{out\_proj} and attention \texttt{o\_proj} groups, which share the
kernel, the grid and the shape, sit 47\,\textmu s apart.

Grouping by name and grid does not help either: the grid $[1,80,1]$ that the two projections
share is also used by the MLP \texttt{down\_proj} of the FP8 layers, which has a different shape.
The split becomes visible only when each launch knows its weight. An Nsight Systems timeline does
carry the information in principle, and a person who already suspects these two projections can
find it by hand; that is how we found it. The point of the per-launch books is that nobody has to
suspect it first.

\begin{table}[t]
\centering\small
\begin{tabular}{@{}lrrrr@{}}
\toprule
Group & launches/step & ms/step & median (\textmu s) & p10--p90 (\textmu s) \\
\midrule
\multicolumn{5}{@{}l}{\emph{By kernel name}} \\
CUTLASS sm120 FP8 GEMM & 148 & 66.67 & 328.3 & 150.0--406.5 \\
\midrule
\multicolumn{5}{@{}l}{\emph{Same launches, by weight, shape and grid}} \\
\texttt{lm\_head} $[5120,248320]$ & 4 & 22.28 & 5562.2 & 5552.1--5589.1 \\
\texttt{linear\_attn.in\_proj\_qkvz} $[5120,16384]$ & 48 & 17.98 & 374.2 & 368.9--379.8 \\
\texttt{linear\_attn.out\_proj} $[6144,5120]$ & 48 & 9.33 & \textbf{194.6} & 190.6--197.5 \\
\texttt{mlp.gate\_up\_proj} $[5120,34816]$ & 8 & 6.25 & 779.7 & 774.2--791.7 \\
\texttt{self\_attn.qkv\_proj} $[5120,14336]$ & 16 & 5.24 & 327.2 & 321.6--332.5 \\
\texttt{mlp.down\_proj} $[17408,5120]$ & 8 & 3.23 & 404.3 & 395.7--410.5 \\
\texttt{self\_attn.o\_proj} $[6144,5120]$ & 16 & 2.36 & \textbf{147.3} & 144.5--149.6 \\

\bottomrule
\end{tabular}
\caption{One CUTLASS kernel, grouped by name and by what each launch read. Batch 1, 40 steps.}\label{tab:namelevel}
\end{table}

\subsection{Capture overhead}\label{sec:overhead}

We ran the batch-1 decode window six times at a locked 2000\,MHz clock, alternating runs without
and with the capture library preloaded. Without it, steps took 122.06, 122.24 and 122.63\,ms;
with it, 122.64, 122.91 and 122.52\,ms. The mean difference is 0.38\,ms (0.3\%), smaller than the
spread among the runs without the library (0.57\,ms), and the last pair is inverted. Outside the
window the library records only graph nodes at instantiation, which happens during warm-up. The
uninstrumented baseline here has no CUPTI subscriber at all, so the figure includes the cost of
kernel activity records, which any tracing profiler pays as well.

\subsection{Coverage}\label{sec:coverage}

Of the 972 launches per step in the batch-1 window, 927 are graph kernel nodes and 45 are direct
launches. All graph nodes resolve. Of the direct launches, 14 per step (560 in the window) carry
no parameters: for these the callback receives neither a function handle nor a symbol name,
\texttt{cudaGetFuncBySymbol} fails, and the host stub has no symbol \texttt{dladdr} can find, so
neither the runtime lookup nor the by-name fallback applies. We suspect the launching code uses a
cudart other than the one the lookup goes through, but have not confirmed it. All 14 are small
PyTorch elementwise, fill and indexing kernels; none streams a weight. \sys{} attributes 25.47\,GB
of weight reads per step, including all four \texttt{lm\_head} reads, which matches the 25.5\,GB
we had computed by hand, before \sys{} existed, from the weight sizes and the number of forward
passes per step.

An earlier version of the analysis reported 22 launches per step without parameters, among them
one \texttt{lm\_head} read. A runtime launch reaches both the runtime and the driver callback under
the same correlation id, and the driver record can come without parameters; the analysis kept
whichever record came last. It now keeps the one that carries parameters.

\subsection{Hardware counters}\label{sec:ncu}

The timing experiments of \cref{sec:split} infer the write-back from durations. To observe it
directly we profiled the isolated GEMM with Nsight Compute, in application replay mode and with
Nsight Compute's own cache flushing disabled, so that each measured launch sees the L2 state its
predecessor left. Before every GEMM a 256\,MB read sweep empties L2 without leaving dirty lines;
then a predecessor either reads or writes $N$\,MiB. On GB10, Nsight Compute reports the \texttt{dram\_\_bytes} metrics as not available; we use L2 counters instead (\cref{tab:ncu}).

\begin{table}[t]
\centering\small
\begin{tabular}{@{}lrrrr@{}}
\toprule
Predecessor & GEMM time & L2 requests & L2 fills from memory & L2 data-bank accesses \\
 & (\textmu s) & (sectors) & (sectors) & (sectors) \\
\midrule
none & 154--160 & 1\,250\,k & 986\,k & 2\,246\,k \\
read 4\,MiB & 169 & 1\,250\,k & 986\,k & 2\,265\,k \\
write 4\,MiB & 171 & 1\,250\,k & 986\,k & 2\,352\,k \\
read 12\,MiB & 158 & 1\,250\,k & 986\,k & 2\,244\,k \\
write 12\,MiB & 198 & 1\,250\,k & 986\,k & 2\,510\,k \\
read 24\,MiB & 153 & 1\,250\,k & 986\,k & 2\,244\,k \\
write 24\,MiB & 213 & 1\,250\,k & 986\,k & 2\,595\,k \\
\bottomrule
\end{tabular}
\caption{Nsight Compute on the isolated $6144\times5120$ FP8 GEMM at $M=4$ after a read-only or a
writing predecessor (medians of three launches; \texttt{lts\_\_t\_sectors}, \texttt{lts\_\_d\_sectors\_fill\_sysmem}, \texttt{lts\_\_d\_sectors}).}\label{tab:ncu}
\end{table}

The GEMM's own traffic is the same in every case to within 0.2\%: 1.25\,M sectors requested and
986\,k sectors (31.6\,MB; the weight is 31.5\,MB) filled from memory. A read-only predecessor
leaves the counters unchanged. A writing predecessor adds L2 data-bank accesses during the GEMM that grow
with the bytes written: 107\,k, 264\,k and 349\,k sectors more than without a predecessor for
4, 12 and 24\,MiB, while the GEMM
takes 171, 198 and 213\,\textmu s. Under replay the timings are noisier than in
\cref{sec:split} (the runs without a predecessor span 154--160\,\textmu s, and the 4\,MiB read case
took 169\,\textmu s with almost no extra data-bank traffic), so the counters are the cleaner
signal here. We read the extra accesses as dirty lines being read out of
the data banks for write-back; the counter does not name its cause, but nothing else in the
launch differs. They amount to 81\% and 67\% of the sectors written for 4 and 12\,MiB; the rest is presumably
written back before the GEMM starts. At 24\,MiB, a full L2, the share drops to 44\%.

\subsection{Dead writes without knowing the model}\label{sec:deadwrites}

The access tracer was not available when we found the GDN stores. Run afterwards, it reaches the
same place without being told where to look. We traced a replay of the vLLM 0.31.0 recording of the
first case (the same kernels at the same addresses, batch 1, every graph launched node by node):
all 974 launches per step are seen (two more than in the engine capture of
\cref{sec:case}; we did not track down the difference), the 269 tensor-map launches included. Four traces of the same
prompt with different seeds and temperatures give 44 resolved steps (\cref{tab:dead}).

\begin{table}[t]
\centering\small
\begin{tabular}{@{}lrrrrrr@{}}
\toprule
Kernel & launches & written & dead & fixed & moving & varies \\
\midrule
GDN decode (\texttt{gdn\_decode\_post\_conv\_mtp}) & 48 & 606.34 & 452.98 & 0.00 & 452.98 & 0.00 \\
causal conv1d update & 48 & 9.83 & 2.95 & 0.00 & 2.95 & 0.00 \\
zero fill (\texttt{triton\_poi\_fused\_zeros}) & 47 & 2.31 & 2.31 & 2.31 & 0.00 & 0.00 \\
temperature & 1 & 3.97 & 0.79 & 0.00 & 0.00 & 0.79 \\
slot mappings & 3 & 0.79 & 0.79 & 0.79 & 0.00 & 0.00 \\
FP4 activation quantization & 72 & 6.68 & 0.28 & 0.28 & 0.00 & 0.00 \\
\midrule
All launches & 974 & 711.87 & 460.18 & 3.43 & 455.96 & 0.79 \\
\bottomrule
\end{tabular}
\caption{Dead writes per decode step of Qwen3.8-27B (MB per step, mean over 44 steps of four traces).}\label{tab:dead}
\end{table}

Of 712\,MB written per step, 460\,MB is overwritten before anyone reads it, and 453\,MB of that
comes from the GDN decode kernel: three of the four per-position states, as the source reading in
\cref{sec:split} predicted. The analysis classifies all of it as moving. As many bytes are dead in
every step, but which ones depends on how many drafts were accepted, which is not known when they
are written; they cannot be dropped, only deferred, which is the fix of \cref{sec:case}. The
causal-convolution state behaves the same way. Only 3.4\,MB per step is fixed, mostly zero fills
that nobody reads, and 0.8\,MB varies: the temperature kernel scales the logits of all four
positions, and the last is read only when every draft is accepted.

Dead in the dataflow sense means never read. A second, value-level test asks whether the output
depends on a write at all: a replay harness poisons one launch's output tensor right after the
launch and compares the generated tokens with a baseline. Of 3\,172 such experiments, 934, covering 476.3\,MB of a step's writes, left the tokens unchanged
over three steps. Rerun over 24 steps, 853 of them (473.8\,MB) still did; 47 (1.3\,MB) changed the
tokens only after the third step, mostly the FP4-quantized activations, and 34 (1.2\,MB) disagreed
between repetitions for reasons we have not found, so a verdict is good to about 1\,MB. In the
traced request the first step accepted every draft, so only the last GDN position was read and the
other three stayed unused over all 24 steps.

\subsection{Cost of access tracing}\label{sec:tracecost}

We have not measured the access tracer's overhead systematically. As an order of magnitude, one
batch-1 decode step of Qwen3-30B-A3B with CUDA graphs took 20.6\,ms without the tracer and about
1.7\,s with it, roughly 80 times slower, including a synchronization and a tensor comparison around
every tensor-map launch. The tracer is meant for a few steps once the per-launch books point at
something, not for continuous use.

